% Chapter 5 — tell the story: Origin and Spread become dated events (kym_events).

\gmenter{events}
\note{Chapter 5: what happened, when, and where. Three stops: Doge's story as a timeline, how
one sentence becomes one event, and the whole corpus.}

\gmgo{events-doge}{Doge's story}
\begin{frame}{Doge, step 4: what happened, when, where?}
  \centering
  \resizebox{!}{6.0cm}{% generated by scripts/figures.py — do not edit
\begin{tikzpicture}[x=1cm, y=1cm, ev/.style={circle, draw=gm@col@events, line width=.8pt, minimum size=3.2mm, inner sep=0pt}]
\node[ev, fill=gm@col@events] (e0) at (0.00,0.00) {};
\node[font=\fontsize{6}{7}\selectfont\bfseries, text=gmInk, above=.9mm of e0] {24 Jun 2005};
\node[font=\fontsize{5.6}{6.6}\selectfont, text=gmInk2, align=center, below=.9mm of e0] {Homestar Runner\\says ``d-o-g-e''};
\node[ev, fill=gm@col@events] (e1) at (1.90,0.00) {};
\node[font=\fontsize{6}{7}\selectfont\bfseries, text=gmInk, above=.9mm of e1] {28 Oct 2010};
\node[font=\fontsize{5.6}{6.6}\selectfont, text=gmInk2, align=center, below=.9mm of e1] {/r/Ads: a photo\\``\dots THIS DOGE''};
\node[ev, fill=gm@col@events!55!gmSurface] (e2) at (3.80,0.00) {};
\node[font=\fontsize{6}{7}\selectfont\bfseries, text=gmInk, above=.9mm of e2] {Apr 2012};
\node[font=\fontsize{5.6}{6.6}\selectfont, text=gmInk2, align=center, below=.9mm of e2] {Tumblr: an audio\\adventure game};
\node[ev, fill=gmSurface, dashed] (e3) at (5.70,0.00) {};
\node[font=\fontsize{6}{7}\selectfont\bfseries, text=gmInk, above=.9mm of e3] {no date};
\node[font=\fontsize{5.6}{6.6}\selectfont, text=gmInk2, align=center, below=.9mm of e3] {the blog\\Your Daily Doge};
\node[ev, fill=gm@col@events] (e4) at (7.60,0.00) {};
\node[font=\fontsize{6}{7}\selectfont\bfseries, text=gmInk, above=.9mm of e4] {7 May 2012};
\node[font=\fontsize{5.6}{6.6}\selectfont, text=gmInk2, align=center, below=.9mm of e4] {YouTube: a fake\\Pok\'emon battle};
\node[ev, fill=gmSurface, dashed] (e5) at (9.50,0.00) {};
\node[font=\fontsize{6}{7}\selectfont\bfseries, text=gmInk, above=.9mm of e5] {no date};
\node[font=\fontsize{5.6}{6.6}\selectfont, text=gmInk2, align=center, below=.9mm of e5] {4chan /v/:\\doge threads};
\node[ev, fill=gm@col@events!55!gmSurface] (e6) at (11.40,0.00) {};
\node[font=\fontsize{6}{7}\selectfont\bfseries, text=gmInk, above=.9mm of e6] {May 2012};
\node[font=\fontsize{5.6}{6.6}\selectfont, text=gmInk2, align=center, below=.9mm of e6] {Tumblr:\\``Polite Doge''};
\node[ev, fill=gm@col@events!55!gmSurface] (e7) at (13.30,0.00) {};
\node[font=\fontsize{6}{7}\selectfont\bfseries, text=gmInk, above=.9mm of e7] {Aug 2012};
\node[font=\fontsize{5.6}{6.6}\selectfont, text=gmInk2, align=center, below=.9mm of e7] {the blog\\F--k Yeah Doge};
\node[ev, fill=gm@col@events!25!gmSurface] (e8) at (13.30,-2.25) {};
\node[font=\fontsize{6}{7}\selectfont\bfseries, text=gmInk, above=.9mm of e8] {2012};
\node[font=\fontsize{5.6}{6.6}\selectfont, text=gmInk2, align=center, below=.9mm of e8] {Shiba Confessions:\\``shibes''};
\node[ev, fill=gm@col@events!55!gmSurface] (e9) at (11.40,-2.25) {};
\node[font=\fontsize{6}{7}\selectfont\bfseries, text=gmInk, above=.9mm of e9] {Dec 2012};
\node[font=\fontsize{5.6}{6.6}\selectfont, text=gmInk2, align=center, below=.9mm of e9] {Reddit:\\/r/DogsIWannaHug};
\node[ev, fill=gm@col@events!55!gmSurface] (e10) at (9.50,-2.25) {};
\node[font=\fontsize{6}{7}\selectfont\bfseries, text=gmInk, above=.9mm of e10] {Dec 2012};
\node[font=\fontsize{5.6}{6.6}\selectfont, text=gmInk2, align=center, below=.9mm of e10] {Cheezburger:\\``Schnauze''};
\node[ev, fill=gm@col@events] (e11) at (7.60,-2.25) {};
\node[font=\fontsize{6}{7}\selectfont\bfseries, text=gmInk, above=.9mm of e11] {8 Jan 2013};
\node[font=\fontsize{5.6}{6.6}\selectfont, text=gmInk2, align=center, below=.9mm of e11] {Reddit: /r/Doge\\is created};
\node[ev, fill=gm@col@events!55!gmSurface] (e12) at (5.70,-2.25) {};
\node[font=\fontsize{6}{7}\selectfont\bfseries, text=gmInk, above=.9mm of e12] {May 2013};
\node[font=\fontsize{5.6}{6.6}\selectfont, text=gmInk2, align=center, below=.9mm of e12] {Reddit:\\/r/dailydoge};
\node[ev, fill=gm@col@events!55!gmSurface] (e13) at (3.80,-2.25) {};
\node[font=\fontsize{6}{7}\selectfont\bfseries, text=gmInk, above=.9mm of e13] {Jul 2013};
\node[font=\fontsize{5.6}{6.6}\selectfont, text=gmInk2, align=center, below=.9mm of e13] {the blog\\shibe-doge};
\node[ev, fill=gm@col@events] (e14) at (1.90,-2.25) {};
\node[font=\fontsize{6}{7}\selectfont\bfseries, text=gmInk, above=.9mm of e14] {29 Jul 2013};
\node[font=\fontsize{5.6}{6.6}\selectfont, text=gmInk2, align=center, below=.9mm of e14] {4chan /s4s/:\\600 replies};
\node[ev, fill=gm@col@events] (e15) at (0.00,-2.25) {};
\node[font=\fontsize{6}{7}\selectfont\bfseries, text=gmInk, above=.9mm of e15] {20 Nov 2013};
\node[font=\fontsize{5.6}{6.6}\selectfont, text=gmInk2, align=center, below=.9mm of e15] {YouTube's Comic\\Sans easter egg};
\draw[gm@col@events, line width=.7pt, -{Stealth[length=1.6mm]}] (e0) -- (e1);
\draw[gm@col@events, line width=.7pt, -{Stealth[length=1.6mm]}] (e1) -- (e2);
\draw[gm@col@events, line width=.7pt, -{Stealth[length=1.6mm]}] (e2) -- (e3);
\draw[gm@col@events, line width=.7pt, -{Stealth[length=1.6mm]}] (e3) -- (e4);
\draw[gm@col@events, line width=.7pt, -{Stealth[length=1.6mm]}] (e4) -- (e5);
\draw[gm@col@events, line width=.7pt, -{Stealth[length=1.6mm]}] (e5) -- (e6);
\draw[gm@col@events, line width=.7pt, -{Stealth[length=1.6mm]}] (e6) -- (e7);
\draw[gm@col@events, line width=.7pt, -{Stealth[length=1.6mm]}] (e7.east) to[out=0, in=0, looseness=1.6] (e8.east);
\draw[gm@col@events, line width=.7pt, -{Stealth[length=1.6mm]}] (e8) -- (e9);
\draw[gm@col@events, line width=.7pt, -{Stealth[length=1.6mm]}] (e9) -- (e10);
\draw[gm@col@events, line width=.7pt, -{Stealth[length=1.6mm]}] (e10) -- (e11);
\draw[gm@col@events, line width=.7pt, -{Stealth[length=1.6mm]}] (e11) -- (e12);
\draw[gm@col@events, line width=.7pt, -{Stealth[length=1.6mm]}] (e12) -- (e13);
\draw[gm@col@events, line width=.7pt, -{Stealth[length=1.6mm]}] (e13) -- (e14);
\draw[gm@col@events, line width=.7pt, -{Stealth[length=1.6mm]}] (e14) -- (e15);
\end{tikzpicture}
}\par\vskip2mm
  {\normalsize\color{gmInk2}
  \tikz\node[circle, draw=gm@col@events, fill=gm@col@events, minimum size=2.6mm, inner sep=0pt]{}; to the day\quad
  \tikz\node[circle, draw=gm@col@events, fill=gm@col@events!55!gmSurface, minimum size=2.6mm, inner sep=0pt]{}; to the month\quad
  \tikz\node[circle, draw=gm@col@events, fill=gm@col@events!25!gmSurface, minimum size=2.6mm, inner sep=0pt]{}; to the year\quad
  \tikz\node[circle, draw=gm@col@events, dashed, fill=gmSurface, minimum size=2.6mm, inner sep=0pt]{}; no date
  \quad---\quad page order}
  \note{
  \begin{itemize}
    \item Layman version: the page tells a story in sentences. We turn it into a timeline: each
      event with its date, its place, its people.
    \item Doge's story, read like a snake: 2005, a cartoon (Homestar Runner) misspells ``dog'' as
      ``doge''. 2010, a Reddit post. 2012, Tumblr, YouTube, 4chan, the first fan blogs. 2013,
      Reddit communities, a 4chan thread with 600 replies, and YouTube's Comic Sans easter egg.
    \item 16 events, in the order the page tells them. A full dot is a date to the day; lighter
      dots are only a month or a year; a dashed one has no date --- and we say so rather than
      guess.
  \end{itemize}}
\end{frame}

\gmgo{events-sentence}{One sentence, one event}
\begin{frame}{One sentence, one event --- every sentence in one}
  \vskip1mm
  \begin{tikzpicture}
    \node[fill=gmTint, rounded corners=2mm, inner sep=3mm, text width=13.4cm, align=left] {\large
      ``\colorbox{gm@col@events!22!gmSurface}{On May 7th}, YouTuber
      \colorbox{gm@col@events!22!gmSurface}{KwandaoRen66} uploaded a video with a person reading the
      text over a fake Pok\'emon battle.''\par\vskip1mm
      {\normalsize\color{gmMuted}Doge, Spread --- the sentence before says ``Sometime in April 2012''}};
  \end{tikzpicture}
  \vskip3mm
  \begin{tikzpicture}[x=1cm, y=1cm,
      fld/.style={rounded corners=2mm, fill=gm@col@events, text=white, minimum width=3.15cm, minimum height=.95cm,
                  align=center, font=\large\bfseries}]
    \foreach \i/\k/\v in {0/when/{7 May 2012}, 1/where/{YouTube}, 2/who/{KwandaoRen66}, 3/certainty/{confirmed}} {
      \node[fld] (f\i) at (\i*3.45,0) {\v};
      \node[above=.3mm of f\i, font=\small, text=gmInk2] {\k};}
  \end{tikzpicture}
  \vskip3mm
  \begin{tikzpicture}[x=1cm, y=1cm,
      st/.style={circle, minimum size=.85cm, fill=gm@col@events!14!gmSurface, draw=gm@col@events, line width=1pt,
                 text=gm@col@events, font=\fontsize{13}{13}\selectfont, inner sep=0pt},
      tx/.style={text width=3.2cm, align=center, font=\small, text=gmInk},
      go/.style={-{Stealth[length=2.4mm]}, line width=1.2pt, draw=gmBase, shorten <=1mm, shorten >=1mm}]
    \foreach \i/\ic/\t in {0/\faRobot/{a model reads the numbered sentences},
                           1/\faHandPointRight/{it points at sentences and names when, where, who},
                           2/\faCalendarCheck/{the pipeline copies the words and reads the date},
                           3/\faClipboardCheck/{an audit: every sentence in an event}} {
      \node[st] (s\i) at (\i*3.45,0) {\ic};
      \node[tx, below=1mm of s\i] {\t};}
    \foreach \i [count=\j] in {0,1,2} \draw[go] (s\i) -- (s\j);
  \end{tikzpicture}
  \note{
  \begin{itemize}
    \item One event, read by a language model --- ministral-3:14b, on the lab's own GPU server ---
      one call per Origin or Spread section.
    \item The model never writes a date or a quote: it points at numbered sentences and returns
      the date's words; the pipeline copies the sentence from the page and parses the date. ``On
      May 7th'' has no year, so the year comes from the sentence before (April 2012); precision
      ``day''. ``Sometime in April 2012'' would be ``month''.
    \item Where: YouTube, a platform. Who: the uploader. Certainty: confirmed --- KYM also writes
      ``allegedly'', ``debunked'', and those are kept as such.
    \item The answer must fit a schema, checked before it is stored, and an audit refuses any
      record that leaves a sentence out. On the full run, 15.3\% of sentences had ended up in no
      event --- the model skipped background and reception. Extraction 4.0.0 makes it an
      invariant: each sentence starts an event or continues the one before. Sentences covered:
      84.9\% $\to$ 100\% on the review sections, 83.6\% $\to$ 100\% on a holdout never used for
      tuning.
    \item The model is drawn from EventKG, not copied: SEM's what/when/where/who and EventKG's rule
      that every fact stays traceable to its source --- every event keeps the sentences it came
      from.
  \end{itemize}}
\end{frame}

\gmgo{events-corpus}{All the stories}
\begin{frame}{\NEvents\ events from \NSections\ stories}
  \begin{columns}[t]
    \column{.6\textwidth}
      {\normalsize\color{gmInk2}events per year}\par\vskip1mm
      % generated by scripts/figures.py — do not edit
\begin{tikzpicture}
\begin{axis}[gm axis, ybar, bar width=4.2pt, width=\linewidth, height=4.6cm,
  xmin=1994.2, xmax=2026.8, ymin=0, ymax=10383, axis x line*=bottom, axis y line=none,
  xtick={1995,2000,2005,2010,2015,2020,2025}, x tick label style={/pgf/number format/1000 sep={}},
  every axis plot/.append style={draw=none}, clip=false]
\fill[gmHair, opacity=.55] (axis cs:2022.5,0) rectangle (axis cs:2026.5,10201);
\node[font=\tiny, text=gmInk2, anchor=north east, align=right] at (axis cs:2026.4,10019) {KYM writes up\\a meme later};
\addplot[fill=gm@col@events] coordinates {(1995,87) (1996,129) (1997,100) (1998,153) (1999,228) (2000,237) (2001,260) (2002,321) (2003,470) (2004,552) (2005,748) (2006,1173) (2007,1463) (2008,1638) (2009,1928) (2010,2636) (2011,3547) (2012,3904) (2013,3944) (2014,3990) (2015,4319) (2016,5474) (2017,6483) (2018,6469) (2019,8547) (2020,9108) (2021,8794) (2022,8976) (2023,7839) (2024,5760) (2025,3947) (2026,1965)};
\node[font=\tiny, text=gmInk, anchor=south] at (axis cs:2020,9108) {9,108};
\end{axis}
\end{tikzpicture}

    \column{.37\textwidth}
      \gmbig{\NFramesEventsPct\%}{of entries have a story --- \NEventsPerFrame\ events each}\par\vskip4mm
      {\normalsize\color{gmInk2}how precisely they are dated}\par\vskip1mm
      % generated by scripts/figures.py — do not edit
\begin{tikzpicture}
\begin{axis}[gm hbar, scale only axis, width=.72\linewidth, height=2.6cm, xmax=132678, symbolic y coords={{to the day},{to the month},{to the year},{no date}},
  point meta=explicit symbolic, nodes near coords={\pgfplotspointmeta}]
\addplot[fill=gm@col@events] coordinates {(91502,{to the day}) [64\%] (7854,{to the month}) [6\%] (6903,{to the year}) [5\%] (36528,{no date}) [26\%]};
\end{axis}
\end{tikzpicture}

  \end{columns}
  \note{
  \begin{itemize}
    \item \NSections\ Origin and Spread sections read; \NEvents\ events; \NFramesEvents\ entries have at
      least one (\NFramesEventsPct\% --- almost every entry that has a story), \NEventsPerFrame\ events each
      on average.
    \item Dates: \NDayPct\% to the day, \NMonthPct\% to the month, \NYearPct\% to the year, \NUndatedPct\% undated. \NPlacePct\% of events
      happen on a platform (YouTube, Twitter, TikTok \dots).
    \item The chart: the history of memes as KYM tells it. The right end falls because KYM writes
      a meme up after it becomes notable --- recent years are not finished yet.
    \item Page order is not time order: KYM tells flashbacks on purpose. The chain between events
      follows the page; the dates carry the time.
  \end{itemize}}
\end{frame}
